\documentclass[11pt]{article}
\usepackage[T1]{fontenc}
\usepackage{textcomp}
\usepackage[margin=0.75in]{geometry}
\usepackage{amsmath,amssymb,amsthm}
\usepackage{array,booktabs}
\usepackage{hyperref}
\setlength{\parindent}{0pt}
\newtheorem{theorem}{Theorem}
\theoremstyle{definition}
\newtheorem{definition}{Definition}
\title{Flow Proof Artifact: isosceles-equal-angles-imply-equal-sides}
\author{Generated by \texttt{flow doc proof}}
\date{Flow Proof Book}
\begin{document}
\maketitle
Equal base angles imply equal opposite sides.\\[0.5em]
\textbf{Source.} euclid — Elements, Book I, Proposition 6\\[0.5em]
\bigskip
\noindent{\large \textbf{Derived fact 1.} \textit{if angle at B equals angle at C then side AB equals side AC} \hfill the Euclidean plane \cdot isosceles triangle \cdot equal base angles imply equal sides}\par
\smallskip\noindent\textit{Coordinate: the Euclidean plane \cdot isosceles triangle \cdot equal base angles imply equal sides}\par
\smallskip\noindent\textit{Source: euclid}\par
\smallskip\noindent\textit{Built on: the angles at the base of an isosceles triangle are equal, two sides and the included angle determine a triangle up to congruence}\par
\medskip
\noindent\textbf{Goal.} if angle at B equals angle at C then side AB equals side AC\par
\noindent$\displaystyle \text{side AB} = \text{side AC}$\par
\medskip
\noindent
\begin{tabular}{@{} >{\raggedright\arraybackslash}p{0.06\textwidth} >{\raggedright\arraybackslash}p{0.47\textwidth} >{\raggedleft\arraybackslash}p{0.41\textwidth} @{}}
\textbf{\#} & \textbf{Proof} & \textbf{Mathematics} \\
\midrule
\textbf{1.} & We prove that equal base angles imply equal sides for isosceles triangles in the Euclidean plane. & \\[0.45em]
\textbf{2.} & We invoke the derived fact governing isosceles triangles in the Euclidean plane: the angles at the base of an isosceles triangle are equal. & \\[0.45em]
\textbf{3.} & We invoke the axiom governing triangle congruence in the Euclidean plane: two sides and the included angle determine a triangle up to congruence. & \\[0.45em]
\textbf{4.} & From step 2 and step 3, this implies side AB equals side AC. Hence proven. & $\displaystyle \text{side AB} = \text{side AC}$ \\[0.45em]
\bottomrule
\end{tabular}
\label{thm:Geometry-isosceles triangle-equal_base_angles_imply_equal_sides}
\par\medskip\hrule
\end{document}
